\documentclass[11pt]{article}

\newcommand{\cnum}{Math 245B}
\newcommand{\ced}{Winter 2026}
\newcommand{\ctitle}[4]{\title{\vspace{-0.5in}\cnum, \ced\\Assignment #1: #2}\author{\vspace{-0.35in}\\#3, #4}}
\usepackage{enumitem}
\usepackage[usenames,dvipsnames,svgnames,table,hyperref]{xcolor}
\usepackage{amsmath, amsfonts}
\usepackage{graphicx} % Required for inserting images
\usepackage{float}
\usepackage{listings}
\usepackage{xcolor}
\usepackage{hyperref}
\usepackage{comment}

\renewcommand*{\theenumi}{\alph{enumi}}
\renewcommand*\labelenumi{(\theenumi)}
\renewcommand*{\theenumii}{\roman{enumii}}
\renewcommand*\labelenumii{\theenumii.}
\newcommand{\notiff}{%
  \mathrel{{\ooalign{\hidewidth$\not\phantom{"}$\hidewidth\cr$\iff$}}}}

\definecolor{codegreen}{rgb}{0,0.6,0}
\definecolor{codegray}{rgb}{0.5,0.5,0.5}
\definecolor{codepurple}{rgb}{0.58,0,0.82}
\definecolor{backcolour}{rgb}{0.95,0.95,0.92}
\definecolor{header}{HTML}{205459}
\definecolor{solution}{HTML}{204d52}

\newcommand{\solution}[1]{{{\textcolor{header}{Solution:} \textcolor{solution}{#1}}}}
\setlist[enumerate]{label=(\alph*)}

\lstdefinestyle{mystyle}{
    backgroundcolor=\color{backcolour},
    commentstyle=\color{codegreen},
    keywordstyle=\color{magenta},
    numberstyle=\tiny\color{codegray},
    stringstyle=\color{codepurple},
    basicstyle=\ttfamily\footnotesize,
    breakatwhitespace=false,
    breaklines=true,
    captionpos=b,
    keepspaces=true,
    numbers=left,
    numbersep=5pt,
    showspaces=false,
    showstringspaces=false,
    showtabs=false,
    tabsize=2
}


\begin{document}
\ctitle{\#1}{}{Asher Christian}{006-150-286}
\date{}
\maketitle

\section{monday}
Lemma: $\{X_n\}$ is sub-martingale and  $T$ is a stopping time and  $T \le n$ almost sure,
then
 \[
\mathbb{E}X_0 \le \mathbb{E}X_t \le \mathbb{E}X_n
\] 
Lemma: $\{X_n\}$ is non-negative sub martingale. 
$X_n^{*} := \max_{0 \le k \le n} X_k$
for all $\lambda \ge 0$
\[
    \lambda \mathbb{P}(X_n^{*} \ge \lambda) \le \mathbb{E}(X_n 1_{X_n^{*} \ge \lambda})
\] 
bound the largest with the last one. Eventually
\[
\mathbb{E}(X_n^{*})^{p} \le \frac{p}{p-1} \mathbb{E}|X_n|^{p}
\] 
\[
\mathbb{E}X^{p} \le \int_{}^{}py^{p-1} \mathbb{P} (X > y)dy
\] 
Proof of first claim:
Let $T := \min_k\{1 \le k \le n ,n : X_k \ge \lambda\}$  to be the first time $X_k$ greater than or equal to lambda or just  $n$ 
$\min\{k : X_k \ge \lambda) \wedge n$. then
\[
    \lambda \mathbb{P}(X_n^{*} \ge \lambda) \le \mathbb{E}(X_T 1_{X_n^{*} \ge \lambda}) = \mathbb{E}(X_t) - \mathbb{E}(X_T 1_{X_n^{*} < \lambda})
\] 
\[
    = \mathbb{E}(X_t) - \mathbb{E}(X_n 1_{X_n^{*} < \lambda})
\] 
Recall the previous lemma.
\[
\mathbb{E}(X_t) \le \mathbb{E}X_n
\] 
so
\[
    \le \mathbb{E}(X_n) - \mathbb{E}(X_N 1_{X_n^{*} < \lambda}) \le \mathbb{E}(X_n 1_{X_n^{*} \ge \lambda}) 
\] 
Application
\[
\lambda \mathbb{P}(X_n^{*} \ge \lambda) \le \mathbb{E} X_n
\] 
\[
\mathbb{P}(X_n^{*} \ge \lambda) \le \frac{ \mathbb{E}(X_n)}{\lambda}
\] 
Martingale $L_p$ convergence theorem:\\
$\{X_n\}$ is martingale and
 \[
\sup_n \mathbb{E} |X_n|^{p} < \infty
\] 
$p > 1$.
Then $X_\infty$ exists and $X_n \rightarrow X_\infty$ in $L_p$
\[
\mathbb{E}|X_n - X_\infty|^{p} \rightarrow 0
\] 
\[
    \mathbb{E}|X_{\infty}^{*}|^{p} < \infty \quad \mathbb{E}|X_\infty|^{p} < \infty
\] 
\[
    X_{\infty}^{*} = \sup_n |X_n|
\] 
First
\[
    \mathbb{E}|X_\infty^{*}|^{p} = \lim_{n\to \infty} \mathbb{E}|X_n^{*}|^{p} \quad X_n^{*} = \max_{1 \le k \le n} |X_k|
\] 
\[
\le \lim_{n\to \infty}\left(\frac{p}{p-1} \mathbb{E}|X_n|^{p}\right) \le \frac{p}{p-1} \sup_n \mathbb{E}|X_n|^{p}
\] 
so both are finite and the rsult is dominated and the convergence follows.

\section{for p = 1} 
$\{X_n\}$ martingale and  $ \sup_n \mathbb{E}|X_n| < \infty$
Q: $X_n \rightarrow X_\infty$ in $L_1$?
If we have $\{X_n\}$ is uniformly integrable then  $X_n \rightarrow X_\infty$ in $L_1$ 
If there is $\phi$ such that
\[
    \frac{\phi(x)}{x} \rightarrow \infty
\] and
\[
\sup_n \mathbb{E} \phi(|X_n|) < \infty
\] 
then $X_n$ are uniformly integrable

\section{Topic Levy's Martingale}
 \[
     X_n \rightarrow X_\infty \text{ a.s.}
\] 
$\{X_n\}$ is martingalealpha function
 \[
     \mathbb{E}(\chi(G) | \mathcal{F}_n]
\] 
\[
    X_n = \mathbb{E}[X_{m} | \mathcal{F}_n] \quad m > n
\] 
\[
    X_1 = \mathbb{E}[X_m | \mathcal{F}_1]
\] 
Can we use say
\[
    X_n = \mathbb{E}[X_\infty | \mathcal{F}_n] ?
\] 
No

\section{Wednesday}
Q: If $\{X_n\}$ is a martingale and  $X_n \to X_\infty$ almost sure when $X_n = \mathbb{E}[X_\infty | \mathcal{F}_n]$\\
Counter examples: SRW $X_\infty = 0$.
if $X_n - \mathbb{E}[X_\infty | \mathcal{F}_n] \implies \mathbb{E}X_n = \mathbb{E}X_\infty$\\
Topic: Levy's Martingale.
Let $X$ be a random variable and $ \mathbb{E}|X| < \infty$ and $ \mathcal{F}_n \nearrow $ filtration, then
\[
    X_n = \mathbb{E}[X | \mathcal{F}_n]
\] 
is a martingale. Proof
\[
    \mathbb{E}[X_{n+1}| \mathcal{F}_n] = \mathbb{E}[ \mathbb{E}[X | \mathcal{F}_{n+1}] | \mathcal{F}_n] = \mathbb{E}[X | \mathcal{F}_n] = X_n
\] 
\section{Lemma}: As above, $\{X_n\}$ is uniformly integrable.\\
Lemma: If $X$ is random variable on $(\Omega, \mathcal{F}, \mathbb{P})$ and $ \mathbb{E}|X| <\infty$ then $\{X_g = \mathbb{E}[X | g] : g \subset \mathcal{F} \text{ sigma algebra}\}$ is $U.I$\\
Proof of Lemma: first $|X_g| + | \mathbb{E}[X | g] | \le \mathbb{E}[|X| | g]$
so for all $g$
\[
    \mathbb{E}|X_g| \le \mathbb{E}( \mathbb{E}[|X| | g]) \le \mathbb{E}|X|
\] 
So for all $g, K$
 \[
\mathbb{P}(|X_g| \ge k) \le \frac{ \mathbb{E}|X|}{k}
\] 
remember that UI if for all $\epsilon$ there exists $K$ such that $ \mathbb{E}|X_\alpha| 1_{|X_\alpha| \ge k} \le \epsilon$ for all $\alpha$\\
SO far we proved that $ \mathbb{E} 1_{|X_g| \ge k} \le \epsilon$ for all $g$.
Next
\[
    \mathbb{E}|X_g|1_{|X_g| \ge k} \le \mathbb{E}[ \mathbb{E}[|X| | g]  1_{|X_g| \ge k}] = \mathbb{E} |X| 1_{|X_g| \ge k}
\] 
For all $\delta > 0$ there exists $k$ such that  $ \sup_g \mathbb{P}(|X_g| \le k) \le \delta$.
For all $\epsilon$, exists $\delta$ such that $\sup_{A, \mathbb{P}(A) < \delta} \mathbb{E}|X|1_{A} \le \epsilon$.
Therefore Levy martingale uniform integrable.
\section{Theorem Levy martingale}
$\{X-n\}$,  $X_n = \mathbb{E}[X | \mathcal{F}_n]$ and $ \mathbb{E}|X|^{p} < \infty$ for som $p \ge 1$ then
there exists  $X_\infty, X_n \to X_\infty$ almost sure and in $L_p$ and
 \[
     X_\infty = \mathbb{E}[X| \mathcal{F}_\infty] \quad F_\infty = \sigma(\bigcup_{ n=1}^{\infty} \mathcal{F}_n)
\] 
Levy upward martingale theorem.
Proof: $ \mathbb{E}|X_n| \le \mathbb{E}[ \mathbb{E}[|X| | \mathcal{F}_n] ] \le \mathbb{E}|X| < \infty$
so $X_n \to X_\infty$ almost sure by martingale convergence theorem
and $\{X_n\}$ is Uniformly integrable which implies  $X_n \to X_\infty$ in $L_1$. If $p > 1$ we use martingale  $L_p$ convergence theorem and so we have the same idea and  $X_n \to X_\infty$ in $L_p$ 
\[
    \sup_n \mathbb{E}|X_n|^{p} < \infty \implies X_n \to X_\infty \text{ in Lp}
\] 
We have $X_n \in \mathcal{F}_\infty$ and so the limit is still measurable with respect to $ \mathcal{F}_\infty$
Last one: $X_\infty = \mathbb{E}[X | F_\infty]$ we use $\pi-\lambda$
We need to prove for all $A \in \mathcal{F}_\infty$ $ \mathbb{E}[X_\infty 1_A] = \mathbb{E}[X 1_A]$
With $\pi-\lambda$ we need to prove that $\forall k, A \in \mathcal{F}_k$
\[
    \mathbb{E}[X_\infty 1_A] = \mathbb{E}[x 1_A]
\] 
proof: Fix $k, A$
\[
    \mathbb{E}[X 1_A] = \mathbb{E}[X_k 1_A]
\] 
and
\[
    \mathbb{E}[X_{k+1} 1_A] = \mathbb{E}[ \mathbb{E}[X | F_{k+1}] 1_A] = \mathbb{E}[X 1_A]
\] 
so for all $n \ge k$
 \[
     \mathbb{E}[X_n 1_A] = \mathbb{E}[X 1_A]
\] 
Then
\[
    \mathbb{E}[X_\infty 1_A] = \mathbb{E}[ \lim_{n\to \infty} X_n 1_A] = \lim_{n\to \infty} \mathbb{E}[X_n 1_A] = \lim_{n\to \infty } \mathbb{E}[X 1_A] = \mathbb{E}[X 1_A]
\] 
where we can swap limit since $X_n \to X-\infty$ in $L_1$ so
\[
   0 = \lim_{n\to \infty} \mathbb{E}[|X_n - X_\infty|] \ge \lim_{n\to \infty}\mathbb{E}[|X_n 1_A - X_\infty 1_A|] 
\] 
If Levy martingale  $X \to \{X_n\}$ and  $X_\infty$ is the limit and $ X_\infty = \mathbb{E}[X | \mathcal{F}_\infty]$
then $X_n = \mathbb{E}[X_\infty | \mathcal{F}_n]$ now if $X_n \to X_\infty$ when does
\[
    X_n - \mathbb{E}[X_\infty | \mathcal{F}_n]
\] 


\end{document}
